\section{Experiments}
\label{sec:exp}

We ask four questions. Does verdict training overexpose facts (\cref{sec:exp-over})? Do overexposed facts surface when the rule is edited (\cref{sec:exp-edit})? What does the repair cost (\cref{sec:exp-cost})? And what else does a state reader buy (\cref{sec:exp-more})? Every hypothesis below was registered before the corresponding run; we report all of them, including those that failed (\cref{app:tables}).

\paragraph{Datasets.} \textbf{InspecSafe-V1}~\cite{inspecsafe}, our primary dataset: 5,013 images from 41 inspection robots at 2,239 waypoints in five industries (tunnel, coal conveyor, oil \& gas, power, metallurgy), 15 hazard facts, a real severity table per industry and the ordinal ``most severe hazard'' rule. Facts are extracted from free-text descriptions and are therefore incomplete. Its rule edits R1--R3 are regulation revisions (stricter PPE; relaxed PPE; relaxed phone and smoking rules). The 107 test images whose level they change, and 21 unchanged controls, were relabelled blind by three annotators with an industrial-safety background who did not take part in model development and saw neither the original levels nor any model output. \textbf{BDD-OIA}~\cite{xu2020bddoia}, the real-image benchmark of the reasoning-shortcut literature~\cite{bortolotti2024rsbench}: six stop reasons with independent multi-label annotations and the ``stop if any reason'' rule, which reproduces the stop label on 99.6\% of images; it gives a clean measurement of overexposure. Its edit E1 (traffic lights may be ignored) follows rsbench and is a controlled probe rather than a realistic policy. \textbf{MMAD}~\cite{jiang2025mmad}: 38 object categories with 104 defect types as facts, the ``any defect'' rule and exact fact labels. We found the reason names of the BDD-OIA release permuted for three indices and use names verified by inspection (\cref{app:data}).

\paragraph{Models.} All models start from Qwen3-VL-8B-Instruct~\cite{qwen3vl} with LoRA (rank 16) for two epochs; calibration is fitted on validation images. \emph{Verdict-only}: the \lux reader trained with the verdict term of \cref{eq:obj} alone---the supervision a verdict-trained model gets. \emph{Level-fill}: facts filled in from the level constraint, a common weak-supervision baseline. \emph{Fact supervision}: verdicts plus fact annotations, unlabelled facts treated as absent. \emph{Direct-SFT}: the same VLM fine-tuned to answer the verdict question. \lux: \cref{eq:obj}.

\subsection{Verdict training overexposes facts}
\label{sec:exp-over}

\begin{table}[t]\centering\small
\caption{\textbf{Overexposure.} False-positive rate on \emph{visible} cells (facts on hazard-free images: switching one on would change the verdict) and \emph{clipped} cells (absent facts on images whose verdict is already decided). DiD: difference-in-differences of verdict-only against zero-shot (clipped drift minus visible drift); all $p<0.001$, paired bootstrap over images.}
\label{tab:over}
\resizebox{\columnwidth}{!}{%
\begin{tabular}{llccc}\toprule
Dataset & Model & FP visible & FP clipped & DiD\\\midrule
\multirow{3}{*}{InspecSafe} & zero-shot & 0.009 & 0.139 & \\
 & verdict-only & 0.000 & 0.232 & +0.102\\
 & \lux & 0.001 & \textbf{0.050} & \\\midrule
\multirow{3}{*}{MMAD} & zero-shot & 0.066 & 0.132 & \\
 & verdict-only & 0.081 & 0.363 & +0.216\\
 & \lux & 0.060 & \textbf{0.088} & \\\midrule
\multirow{3}{*}{BDD-OIA} & zero-shot & 0.018 & 0.043 & \\
 & verdict-only & 0.030 & 0.199 & +0.143\\
 & \lux & 0.022 & \textbf{0.021} & \\\bottomrule
\end{tabular}}
\end{table}

\Cref{thm:gap} predicts an asymmetric drift: verdict training should corrupt clipped cells, which carry no information, and leave visible cells alone. \Cref{tab:over} shows exactly this on all three datasets and both rule families: false positives on clipped cells grow 1.7--4.6$\times$, while visible cells stay flat or improve. The effect is graded by exposure. On BDD-OIA the two frequent reasons with the lowest training exposure, traffic sign (0.06) and other obstacles (0.04), drift most (0.03$\to$0.33 and 0.03$\to$0.40), and the two most exposed, car (0.35) and red light (0.33), drift least (0.05$\to$0.17, 0.03$\to$0.08). The drift replicates on a second seed (DiD $+0.107$ on BDD-OIA; nine runs on InspecSafe), and \lux removes it in every run.

\FIG{Figure 3. Distributions of fact probabilities on clipped (top) and visible (bottom) cells for zero-shot / verdict-only / \lux, InspecSafe and BDD-OIA. Verdict-only moves mass to 1 on clipped cells only. Data: results/fig1\_industrial/story.json, results/figdata/fig3\_calibration\_stats.json.}

\subsection{Overexposed facts surface under rule edits}
\label{sec:exp-edit}

\begin{table}[t]\centering\small
\caption{\textbf{Real regulation edits on InspecSafe}, scored against blind human labels; accuracy on the images whose level changes (zero new labels). R1 stricter PPE (n=48, upgrades), R2 relaxed PPE (n=43, downgrades), R3 relaxed phone and smoking rules (n=16). \lux (full) vs.\ Direct-SFT with the new rule in its prompt: $+0.63$ [$+0.51$, $+0.73$].}
\label{tab:inspec-edit}
\resizebox{\columnwidth}{!}{%
\begin{tabular}{lcccc}\toprule
 & all 107 & R1 & R2 & R3\\\midrule
Direct-SFT, old prompt & 0.215 & 0.062 & 0.233 & 0.625\\
Direct-SFT, new rule in prompt & 0.262 & 0.167 & 0.233 & 0.625\\
Rule-conditioned Direct-SFT & 0.449 & 0.417 & 0.349 & 0.812\\
\lux zero-shot & 0.832 & 0.958 & 0.651 & 0.938\\
\lux reader, verdict-only training & 0.850 & 0.958 & 0.721 & 0.875\\
\lux reader, level-fill training & 0.804 & 0.958 & 0.605 & 0.875\\
\lux, seed 0 / 1 & \textbf{0.888 / 0.879} & 0.958 & \textbf{0.767} & \textbf{1.000}\\\midrule
\textit{Direct-SFT + 200 relabelled images} & 0.673 & 0.896 & 0.395 & 0.750\\
\textit{Direct-SFT retrained on relabelled data} & 0.570 & 0.812 & 0.442 & 0.188\\\bottomrule
\end{tabular}}
\end{table}

\begin{table}[t]\centering\small
\caption{\textbf{Controlled edit on BDD-OIA} (independent labels). E1: traffic lights may be ignored; E2: other obstacles may be ignored. Accuracy over all 4,572 test images and on the images whose decision changes. Seed 1 replicates every row (\lux 0.818, verdict-only 0.116, Direct-SFT 0.213 on the E1 affected images).}
\label{tab:bdd-edit}
\resizebox{\columnwidth}{!}{%
\begin{tabular}{lcccc}\toprule
 & \multicolumn{2}{c}{E1} & \multicolumn{2}{c}{E2}\\
 & all & affected (670) & all & affected (86)\\\midrule
zero-shot & 0.675 & 0.452 & 0.714 & 0.500\\
verdict-only & 0.695 & 0.166 & 0.799 & 0.453\\
Direct-SFT, edit in prompt & 0.705 & 0.184 & 0.801 & 0.430\\
fact supervision & 0.803 & 0.816 & 0.805 & 0.779\\
\lux & \textbf{0.799} & \textbf{0.810} & 0.801 & 0.756\\\bottomrule
\end{tabular}}
\end{table}

\Cref{thm:gap} predicts which edits are dangerous: those that relax a fact and thereby expose what it clipped. \Cref{tab:inspec-edit} confirms this on real regulations scored against human labels. A fine-tuned verdict model does not follow a rule it is told about (0.262), and even retraining it on relabelled data reaches only 0.570, because the relabelled descriptions inherit their omissions. \lux, with the table edited and nothing retrained, is correct on 0.888 / 0.879 of the affected images (two seeds). The largest gap between training objectives is on the relaxation R2, where clipped facts surface: 0.767 for \lux against 0.605 for level-fill training. \Cref{fig:teaser} shows one such scene. The human labels agree with the description-derived ones on 88\% of the affected images (weighted $\kappa=0.94$), and \lux scores higher against the human labels than against the derived ones, so the derived labels did not favour it.

On BDD-OIA (\cref{tab:bdd-edit}) the verdict-trained model falls from 0.45 (untrained) to 0.17 on the images whose decision changes under E1, and \lux reaches 0.81. Over all images the gain is $+10.4$ points [$+9.2$, $+11.7$], smaller than on the affected subset: on the 382 frames with a red light \emph{and} a real obstacle, the verdict-trained model keeps stopping more often (0.95 vs.\ 0.77), because it over-reads every reason---the same bias that fails on the affected images.

\subsection{What does the repair cost?}
\label{sec:exp-cost}

\begin{table}[t]\centering\small
\caption{\textbf{Incomplete annotations.} Top: MMAD, defect-type labels kept at random in the training set; $\Delta$AUROC relative to verdict-only training. Bottom: BDD-OIA with half of the reason labels missing.}
\label{tab:dose}
\resizebox{\columnwidth}{!}{%
\begin{tabular}{lccc}\toprule
MMAD, labels kept & $\Delta$AUROC & FP clipped & fact acc.\\\midrule
none (verdict-only) & -- & 0.363 & 0.603\\
100\% & $-0.004$ & 0.088 & 0.756\\
78\% & $-0.008$ & 0.080 & 0.717\\
58\% & $-0.009$ & 0.091 & 0.663\\
37\% & $-0.012$ & 0.119 & 0.572\\
58\%, unlabelled = absent & $-0.049$ & 0.082 & 0.662\\\midrule
BDD-OIA, 50\% kept & $\Delta$AUROC & FP clipped & recall\\\midrule
fact supervision, unlabelled = absent & $-0.016$ & 0.016 & 0.552\\
\lux, cross-fitted $c$ & $-0.009$ & 0.027 & 0.593\\
\lux, in-training $c$ & $-0.005$ & 0.039 & 0.629\\
\lux, true $c$ & $-0.007$ & 0.024 & 0.616\\\bottomrule
\end{tabular}}
\end{table}

With complete fact annotations the repair is free: \lux is non-inferior to verdict-only training in verdict AUROC on MMAD ($-0.004$, 95\% CI $[-0.009,+0.000]$) and BDD-OIA ($-0.005$; seed 1: $-0.001$). With incomplete annotations a cost appears and grows monotonically with the missing rate (\cref{tab:dose}). What matters most is how missing annotations are treated: treating them as negatives costs five times more than modelling them (MMAD, 58\%), and on BDD-OIA with half the labels missing it pushes fact recall down to 0.55, whereas \lux keeps it at 0.59--0.63, close to the 0.63 of complete labels. The estimate of $c_j$ itself is reliable only for frequent facts (car: 0.51 against a true 0.50), but the outcome is insensitive to it. InspecSafe, whose facts come from incomplete descriptions, sits on this curve: \lux trails verdict-only training by 0.8--1.4 points of in-distribution accuracy, and the residual errors are almost entirely decisive facts that are hard to read (\cref{sec:exp-health}).

\subsection{What else does a state reader buy?}
\label{sec:exp-more}

\paragraph{Transfer and new facts.} Leaving one InspecSafe industry out of training, \lux (level-fill) reaches 0.930 pooled accuracy against 0.889 for Direct-SFT ($+4.1$, $p<0.001$). A hazard unseen in training and calibration is detected by adding its question and table column (accuracy on the test images where the new hazard alone decides the level; smoking 0.94 / 0.81 over two seeds vs.\ 0.25 for Direct-SFT trained without it, non-motor vehicle 1.00 vs.\ 0.12).

\paragraph{Efficiency and robustness.} A \lux query takes 0.30\,s against 0.29\,s for asking the verdict and 4.38\,s for a reasoning model (median, Qwen3-VL-8B). Three paraphrased question sets change zero-shot accuracy by at most 1.5 points. Composing the facts with the rule beats handing the same fact probabilities to an LLM ($-4.2$ / $-4.8$ points for 8B / 32B aggregators), describing then deciding with the same VLM ($-19.0$), and open-vocabulary detectors composed with the same rules (0.864--0.881 vs.\ 0.919 zero-shot).

\paragraph{Other backbones.} \TBD{V48: Qwen3-VL 2B / 4B and Gemma 4 E4B on \cref{tab:over,tab:inspec-edit}; Qwen3-VL-32B zero-shot.}

\subsection{State health}
\label{sec:exp-health}

\emph{Verdict leakage}---the $R^2$ of an absent fact's logit on the scene level---is 0.011 for the untrained model, 0.064 after verdict-only training and 0.038--0.041 with \lux: a fact that is not there should not know how dangerous the scene is, which is the conditional independence assumed in \cref{thm:gap}. On InspecSafe, the rule applied to the facts of fully described images reproduces the label in all cases; the in-distribution gap to verdict training lies entirely in 251 multi-hazard images, where every under-call misses a decisive fact that is written in the description.

\FIG{Figure 4. Training dynamics on BDD-OIA (V47): old-rule accuracy, clipped-cell FP and E1 accuracy against images seen, verdict-only vs.\ \lux.}
